\documentclass[11pt]{article}

\usepackage[margin=1in]{geometry}
\usepackage{amsmath}
\usepackage{amssymb}
\usepackage{amsthm}
\usepackage{booktabs}
\usepackage{array}
\usepackage[hidelinks]{hyperref}

\theoremstyle{plain}
\newtheorem{theorem}{Theorem}[section]
\newtheorem{lemma}[theorem]{Lemma}
\newtheorem{proposition}[theorem]{Proposition}
\newtheorem{corollary}[theorem]{Corollary}

\theoremstyle{definition}
\newtheorem{definition}[theorem]{Definition}

\theoremstyle{remark}
\newtheorem{remark}[theorem]{Remark}
\newtheorem*{provenance}{Provenance}

\newcommand{\Lb}{\overline{L}}

\title{The polarized-Frobenius axiom class: an axiom statement with
Davenport--Heilbronn and Epstein witnesses (not a theorem)}
\author{Jay Tyagi}
\date{August 26, 2026}

\begin{document}

\maketitle

\begin{abstract}
This note records an axiom class together with a witness record; it establishes no
theorem. A \emph{polarized Frobenius system} for a Dirichlet series $D(s)$ over a
proper adelic curve (in the sense of Chen--Moriwaki, with nef adelic line bundles in
the Yuan--Zhang calculus) consists of six axioms: a polarized square, a Frobenius
family, a total composition law with primitivity exactly at the primes, effectivity,
a transpose/Rosati-type functional-equation template, and a $D$-attachment clause.
We show that the attachment clause cannot be formalized innocently: under a degree
reading, an Euler product is forced by definition, and an explicit
$\mathbb{P}^1 \times \mathbb{P}^1$ system satisfies axioms (PF1)--(PF4) while
touching no coefficient of any $D$; under a trace reading, counting data does not
factor through composition, and an organizing Lefschetz-tie axiom must be added. On
every branch the exclusion of RH-false examples is definitional. With effectivity as
an axiom, a cheap monoid computation shows that admissible series have von Mangoldt
data supported on prime powers and nonnegative. Exact witnesses are recorded: the
Davenport--Heilbronn series and the Epstein zeta function of $x^2 + 5y^2$ fail
admissibility, with closed-form von Mangoldt values; but so does the RH-true
Dirichlet $L$-function $L(s, \chi_4)$. The class is therefore a positivity filter,
not an Euler-product filter and not an RH filter, and Beurling-type generalized
number systems satisfy all of its operative coefficient-level consequences, so the
class as printed cannot decide the Riemann hypothesis. All numerical witnesses carry
exact closed forms.
\end{abstract}

\section{Introduction and status}\label{sec:intro}

This document is an \textbf{axiom statement plus a witness record}. It is NOT a
theorem, it is NOT ``provable today,'' and it does NOT establish that the
Davenport--Heilbronn and Epstein zeta functions admit no
$\mathbb{F}_1$-model (a claim that would require a formalized notion of
$\mathbb{F}_1$-model and remains open). An earlier ``Lemma~B'' theorem framing of
this material was refuted by two independent adversarial reviews (blind
replication) and the refutation was upheld on independent re-derivation in an
adjudication pass; the review record, and the meaning of the internal terms used
below (``killers,'' ``adjudication,'' ``referee''), are summarized in the
Provenance remark at the end of this note. What survives that process --- per the
adjudication's partial overrule of one kill verdict, via a referee-verified minor
finding --- is exactly this: \textbf{with effectivity added as an axiom, the
step-2 monoid computation under the degree reading is a correct, cheap formal
computation}, so the axiom-format note is a legitimate, nonzero item rather than
pure re-dressing. The content remains quasi-definitional: it is a known
positivity-barrier witness in a new dialect, whose sharper formal home is the
Blomer--Leung monoid converse \cite{BL}.

\medskip
\noindent\textbf{Role.} This note is the \textbf{axiom citation} for any argument
that wants to consume ``multiplicativity as geometry'' --- nothing more. It
supersedes the earlier ``Lemma~B (target statement)'' framing, which is retained
in the program record as the immutable historical version.

\medskip
\noindent\textbf{Binding constraints.} Every mathematical claim below was either
(i) re-derived independently by both adversarial reviews and the adjudication, or
(ii) computed and stored with exact closed forms (the Provenance remark traces
every number to its source). Claims struck in adjudication (theorem status,
``provable today,'' discharge of the $\mathbb{F}_1$-model question, any credit
beyond the axiom-format role) do not appear here except in the negative.

\section{The axiom class (repaired definition)}\label{sec:axioms}

Let $S = (K, (\Omega, \nu), (|\cdot|_\omega))$ be a \textbf{proper adelic curve}
in the Chen--Moriwaki sense \cite{CM} (``proper'' $=$ the product formula holds).
A \textbf{polarized Frobenius system} for a Dirichlet series $D(s)$ over $S$
consists of:

\begin{itemize}
\item \textbf{(PF1) Polarized square.} A fibered square $X$ over $S$ carrying a
nef adelic $\mathbb{Q}$-line bundle $\Lb$ in the Yuan--Zhang adelic-line-bundle
calculus \cite{YZ}, with the two fiber classes $\xi_1, \xi_2$ and a
correspondence composition $\circ$ and transpose $t$.
\item \textbf{(PF2) Frobenius family.} A family of correspondences
$\{\Psi_n\}_{n \ge 1}$ on $X$ with $\Psi_1 = \Delta$ (the identity
correspondence).
\item \textbf{(PF3) Composition law and primitivity.}
$\Psi_m \circ \Psi_n = \Psi_{mn}$ for \textbf{all} $m, n \ge 1$, and $\Psi_p$ is
primitive (not a composition of non-identity members) exactly at the primes $p$.
\item \textbf{(PF4) Effectivity} \emph{(axiom added by the adjudicated repair;
the original definition omitted it and its step~3 is false without it --- general
correspondence classes are differences of effective ones and can have negative
degrees)}. Every $\Psi_n$ is effective.
\item \textbf{(PF5) FE/duality (transpose/Rosati-type)} \emph{(axiom added by the
adjudicated repair; the original definition's ``closed under transpose'' was
never tied to the $\Psi_n$ by any axiom)}. The family is stable under transpose,
and a fixed Rosati-type identity ties $\Psi_n^t$ to $\Psi_n$ through the
polarization $\Lb$, whose numerical shadow is the $s \mapsto 1-s$ (eigenvalue
$\alpha \mapsto q/\alpha$) symmetry of the functional equation of $D$. Template:
in characteristic $p$ the functional equation IS Rosati/Poincar\'e duality,
relating $\Gamma_F^t$ to the dual ``$q/F$''
($\Gamma_F^t \circ \Gamma_F = q \cdot \Delta$ for the purely inseparable
Frobenius $F$ of degree $q$). Over an adelic curve this axiom is stated here at
template level only: \textbf{no candidate substrate currently instantiates it for
any $D$}, and this note makes no claim about how a sharpened transpose axiom
would cut down the model class of \S\ref{sec:hornA} (unadjudicated).
\item \textbf{(PF6) $D$-attachment.} The counting data of $D(s)$ is reproduced by
the system. This clause is the definitional joint, and it cannot be formalized
innocently: \S\ref{sec:fork} shows that every reading makes the
Davenport--Heilbronn/Epstein exclusion \textbf{definitional}, which is why this
document is an axiom statement and not a theorem.
\end{itemize}

\medskip
\noindent\textbf{PF5 status (read together with \S\ref{sec:hornA} and
\S\ref{sec:clauses}).} Two bookkeeping facts about (PF5), stated here so that no
reader has to assemble them: \textbf{(a)} PF5's transpose-stability clause
already fails for the model system of \S\ref{sec:hornA} ---
$\Gamma_{z \mapsto z^n}^t$ has bidegree $(n, 1)$, which is not the bidegree
$(1, m)$ of any $\Gamma_m$ --- so that exhibit certifies vacuity of the
\textbf{(PF1)--(PF4) fragment}, and no transpose-stable instantiation of the full
six-axiom class is currently known for any $D$. This is consistent with PF5's
uninstantiated-template status and changes nothing in the fork of
\S\ref{sec:fork}: Horn~A's steps 1 and 3 do not use the exhibit. \textbf{(b)} In
the \S\ref{sec:clauses} accounting, PF5 is excluded from ``the axiom a DMV world
violates'' because it is an uninstantiated template contributing no operative,
checkable coefficient-level consequence. If its FE \textbf{shadow} were instead
imposed directly on $D$ as an analytic requirement, FE-free Beurling systems
would fail it --- but that would be an analytic condition inserted by hand, not
an axiom any substrate enforces, and the class would still sit below the
degree-one-rigidity bracket of \S\ref{sec:clauses} (no Ramanujan input, no
degree-one-rigidity derivation). The positivity-filter classification and the
relocation order of \S\ref{sec:clauses} are unchanged either way.

\section{Why this is an axiom statement and not a theorem: the definitional
fork}\label{sec:fork}

\subsection{Preliminary: bidegrees multiply under composition; traces do
not}\label{sec:bidegrees}

For correspondences on a square, the bidegrees
$(d_1, d_2) = (Z \cdot \xi_1,\, Z \cdot \xi_2)$ are multiplicative under
composition: $d_i(\Psi' \circ \Psi) = d_i(\Psi') \cdot d_i(\Psi)$ (over a
generic point the fibers compose; verified on graph models, four pairs, in the
adjudication). Intersection numbers against the diagonal are NOT:
$(\Psi' \circ \Psi \cdot \Delta)$ is not a function of $(\Psi' \cdot \Delta)$
and $(\Psi \cdot \Delta)$. Concretely, on $\mathbb{P}^1 \times \mathbb{P}^1$
with $\Gamma_n = \operatorname{graph}(z \mapsto z^n)$:
$(\Gamma_n \cdot \Delta) = n + 1$, verified two ways (numerical classes:
$\Gamma_n \equiv n\xi_1 + \xi_2$, $\Delta \equiv \xi_1 + \xi_2$, product
$= n + 1$; fixed points of $z^n = z$: the $n - 1$ roots of $z^{n-1} = 1$ plus
$0$ and $\infty$). And $n + 1$ is not multiplicative:
$(\Gamma_6 \cdot \Delta) = 7 \ne 3 \cdot 4 =
(\Gamma_2 \cdot \Delta)(\Gamma_3 \cdot \Delta)$.

\subsection{Horn A --- the degree reading}\label{sec:hornA}

Read (PF6) as: \emph{the degrees of the $\Psi_n$ against the two fibers
reproduce the counting data of $D$} (the reading the original commission's own
template sentence mandated). Then:

\begin{enumerate}
\item By \S\ref{sec:bidegrees} the attached degree data is \textbf{completely
multiplicative by definition} --- an Euler product is forced by (PF3) before any
geometry is consulted.
\item \textbf{The class touches no coefficient of any $D$.} The system
$X = \mathbb{P}^1 \times \mathbb{P}^1$ over the adelic curve of $\mathbb{Q}$
with the standard nef bundle $\mathcal{O}(1,1)$ and
$\Psi_n = \Gamma_{z \mapsto z^n}$ satisfies every axiom (PF1)--(PF4):
$\Psi_m \circ \Psi_n = \Psi_{mn}$ exactly, $\Psi_p$ primitive exactly at
primes, every $\Psi_n$ an effective irreducible curve, bidegrees $(1, n)$.
Nothing in (PF1)--(PF4) references the coefficients of $D$, so the identical
system would serve $D = \zeta$, $D = \mathrm{DH}$, or $D = \mathrm{Epstein}$
indifferently. (Both adversarial reviews constructed this counterexample
independently; the adjudicator re-verified it. The printed (PF5) already
excludes it --- its transpose-stability clause fails, $\Gamma_n^t$ having
bidegree $(n, 1)$, not any $\Gamma_m$ --- so the exhibit certifies vacuity of
the (PF1)--(PF4) fragment; see the PF5 status paragraph, \S\ref{sec:axioms}.
Whether some sharpened transpose axiom would cut down the model class further is
not adjudicated and not claimed.)
\item Therefore, under Horn~A, ``reproduce the counting data of $D$'' must
itself be an \textbf{axiom --- an admissibility condition on $D$}, not a
property the geometry proves. Under that admissibility axiom, prime-power
support and $\Lambda_D \ge 0$ follow by the correct, cheap monoid computation of
\S\ref{sec:admissibility} --- \textbf{quasi-definitionally}: the only geometric
ingredient consumed is that degrees of effective correspondences against a nef
class are $\ge 0$, which is a basic property of the Yuan--Zhang calculus
\cite{YZ}, not a theorem being put to work.
\end{enumerate}

\subsection{Horn B --- the trace reading}\label{sec:hornB}

Read (PF6) as: \emph{the counting data of $D$ is reproduced by intersection
against the diagonal, $(\Psi_n \cdot \Delta)$-type numbers} --- the only
reading matching where counting data actually lives in the char-$p$ template
(\S\ref{sec:charp}). Then the original step~2 (``composition law $+$
primitivity $\Longrightarrow$ prime-power support; formal monoid algebra; no
analysis'') is \textbf{false as a formal claim}: by \S\ref{sec:bidegrees} trace
data does not factor through composition ($(\Gamma_6 \cdot \Delta) = 7 \ne 12$
on $\mathbb{P}^1 \times \mathbb{P}^1$), and no formal argument forces
$(\Psi_6 \cdot \Delta)$-type data to vanish. What is needed is a
\textbf{Lefschetz-tie axiom} (``the counting data of $D$ is trace-organized by
the monoid'') --- in characteristic $p$ that tie is a THEOREM about the
existing cohomology; over $\mathbb{Q}$ it is exactly the missing substrate.
Adding it as an axiom makes the $D$-attachment Euler-organized by definition,
so the Davenport--Heilbronn/Epstein exclusion again follows from the definition
plus the one-line coefficient witnesses of \S\ref{sec:witnesses} ---
definitional once more.

\subsection{Exhaustiveness of the fork}\label{sec:exhaustive}

The primary dichotomy is formal, not empirical: the counting functional
attached by (PF6) either \textbf{factors through composition} --- through
composition-multiplicative invariants (bidegrees and monomials in them;
equivalently the induced action on $H^0/H^2$) --- which is Horn~A; or it
\textbf{does not}. Any non-multiplicative reading whatsoever --- the $H^1$
trace or any other --- breaks step~2's formal monoid algebra and therefore
needs an organizing tie axiom, which makes the exclusion definitional: Horn~B's
analysis applies verbatim to all of them. The $H^0/H^2$-versus-$H^1$ trichotomy
is the template interpretation of this split (in every known template the zero
side IS the $H^1$ trace), not the load-bearing case division. Under Horn~A the
statement is an admissibility definition; under Horn~B it needs a
Lefschetz-type axiom that makes it an admissibility definition. Hence: axiom
statement, on every branch. (Both adversarial reviews, independently; upheld in
adjudication.)

\subsection{The char-\texorpdfstring{$p$}{p} accounting,
corrected}\label{sec:charp}

The commission's template sentence (``the char-$p$ model --- degrees of
graph-of-Frobenius powers on $C \times C$ reproducing point counts --- is the
template'') was \textbf{factually wrong} and is corrected here. On $C \times C$
over $\mathbb{F}_q$, the bidegrees of $\Gamma_{\mathrm{Frob}^k}$ are
$(1, q^k)$ --- these are only the \textbf{main terms} of
$\#C(\mathbb{F}_{q^k}) = q^k + 1 - \sum_i \alpha_i^k$. The zero side
$-\sum_i \alpha_i^k$ lives \textbf{exclusively in the $\Delta$-intersection}
(Lefschetz: $N_k = (\Gamma_{F^k} \cdot \Delta)$); degrees never see it. A
degree-based reading is therefore structurally blind to zeros. Numerically
verified in the adjudication on $E \colon y^2 = x^3 + x + 1$ over
$\mathbb{F}_7$: bidegrees $(1, 7)$, $(1, 49)$, $(1, 343)$ versus counts
$N_1 = 5$, $N_2 = 55$, $N_3 = 380$ (Frobenius trace $a = 3$:
$\alpha + \beta = 3$, $\alpha\beta = 7$).

\section{The scope restriction: total multiplicativity pins degree-one Euler
shape}\label{sec:scope}

(PF3) demands $\Psi_m \circ \Psi_n = \Psi_{mn}$ for \textbf{all} $m, n$ --- not
merely coprime pairs. Two consequences, both stated here as adjudicated
obligations:

\begin{itemize}
\item \textbf{(a) Degree-one Euler shape.} Matched at the level of the
Dirichlet coefficients $a_n$ of $D$, totally multiplicative coefficient data is
exactly the Euler shape
$D(s) = \prod_p (1 - a_p p^{-s})^{-1}$ --- \textbf{degree one}. The axiom
class as printed excludes every higher-degree Euler-intact $L$-function (e.g.\
$\zeta_K$ for a quadratic field $K$ has degree-2 Euler factors and
non-totally-multiplicative coefficients; it satisfies the coefficient-level
consequence class of \S\ref{sec:admissibility} but not the printed degree-one
shape). A relaxed composition law (e.g.\ coprime-only multiplicativity plus
per-prime tower data) is NOT provided by the printed axioms; any future
widening is new design, outside this note.
\item \textbf{(b) Relation to the one-generator function-field monoid.} The
char-$p$ template's Frobenius monoid is generated by ONE element:
$\{\Gamma_{F^k} : k \ge 0\} \cong (\mathbb{N}, +)$, indexed by the exponent
$k$ (the extension degree), with a single ``prime.'' The printed multi-prime
monoid $\{\Psi_n\}$ over $(\mathbb{N}^\times, \cdot)$ --- free commutative on
all rational primes --- is a different object: the commission's template was a
one-prime template misread as a multi-prime one. Total multiplicativity across
distinct primes is exactly the new, restrictive content: at coefficient level
it pins the degree-one Euler shape (a); if instead the degree data were matched
against $\Lambda$-type (prime-power-supported) counting data directly, complete
multiplicativity would contradict vanishing at products of distinct primes
($c_{pq} = c_p \cdot c_q \ne 0$ whenever two primes carry nonzero data) unless
at most one prime carries data --- i.e.\ exactly the one-generator char-$p$
monoid again. (Both adversarial reviews; referee minor finding; upheld.)
\end{itemize}

\section{The admissibility axiom and the exclusion computation (the correct
cheap part)}\label{sec:admissibility}

\begin{definition}[admissible $D$]\label{def:admissible}
$D(s) = \sum_n a_n n^{-s}$ ($a_1 = 1$) is \emph{admissible} for the
polarized-Frobenius axiom class if, writing
$-D'/D(s) = \sum_n \Lambda_D(n)\, n^{-s}$, the counting data satisfies:
\textbf{$\Lambda_D$ is supported on prime powers and $\Lambda_D(n) \ge 0$ for
all $n$} (with the FE/duality clause (PF5) named as a further requirement whose
substrate-level content is currently uninstantiated). In the strict printed
(degree-one) shape, $\Lambda_D(p^k) = a_p^k \cdot \log p$ with $a_p \ge 0$.
\end{definition}

\noindent\textbf{The computation (correct, cheap, and quasi-definitional).}
Under Horn~A plus the admissibility axiom, the attachment functional is fixed
explicitly: \textbf{$a_n := d_2(\Psi_n) = (\Psi_n \cdot \xi_2)$, the
second-fiber degree} --- the char-$p$ analog of $\deg F^k = q^k$. (This is the
multiplicative choice: by \S\ref{sec:bidegrees} the individual bidegree
components $d_i$, and monomials in them, are monoid homomorphisms under
composition, whereas the degree against the polarization itself --- an
$(\xi_1 + \xi_2)$-type pairing --- is NOT: on graphs,
$(1+n)(1+m) \ne 1 + nm$. Everything below runs verbatim for any declared
multiplicative monomial in the bidegrees; nothing below may be read against the
non-multiplicative polarization pairing.) (PF3) then makes
$n \mapsto a_n = d_2(\Psi_n)$ a monoid homomorphism, i.e.\ totally
multiplicative, forcing the degree-one Euler shape (\S\ref{sec:scope}(a));
primitivity places the generators exactly at the primes; hence $-D'/D$ has
$\Lambda_D(p^k) = a_p^k \log p$ and $\Lambda_D(n) = 0$ off prime powers ---
formal monoid algebra, no analysis. (PF4) effectivity plus (PF1) nefness gives
$a_p = d_2(\Psi_p) = (\Psi_p \cdot \xi_2) \ge 0$, because $\xi_2$ is nef and
$\Psi_p$ is effective, hence $\Lambda_D(n) \ge 0$. This is the step-2/step-3
chain the referee verified by writing out the monoid computation, valid
\textbf{once effectivity is an axiom} --- and it is an unpacking of the
definition, not a theorem about $D$: the geometry contributes only ``nef
against effective is $\ge 0$.''

\medskip
\noindent\textbf{What would restore theorem status.} Nothing currently
available: a genuine exclusion theorem would require a substrate on which the
Lefschetz tie of \S\ref{sec:hornB} is \emph{proved} rather than posited ---
i.e.\ the missing substrate itself. Until then the correct citation form is:
``$D$ is inadmissible for the polarized-Frobenius axiom class [this note,
\S\S\ref{sec:admissibility}--\ref{sec:witnesses}],'' never ``$D$ admits no
polarized Frobenius system [theorem].''

\section{The witnesses (exact)}\label{sec:witnesses}

All values below are stored with exact closed forms; sources are traced in the
Provenance remark. ``Inadmissible'' always means: fails the
\S\ref{sec:admissibility} admissibility condition, with the named witness.

\subsection{Davenport--Heilbronn (RH-false, and inadmissible)}\label{sec:dh}

Coefficient system: $a_n$ periodic mod $5$,
$(a_1, \ldots, a_5) = (1, \kappa, -\kappa, -1, 0)$, with the Gauss-sum surd
\begin{equation*}
\kappa = \frac{\sqrt{10 - 2\sqrt{5}} - 2}{\sqrt{5} - 1}
= 0.28407904384041229602829183239312616909109
\end{equation*}
($41$ digits; last digit rounded --- the true continuation is
$\ldots 9091088\ldots$).

DH is entire, satisfies a functional equation (verified to
$5.6 \times 10^{-40}$), and has the verified off-line zero
$\rho = 0.808517182456637 + 85.699348485377592\,i$ (Newton residual
$7.6 \times 10^{-41}$). Its counting data violates BOTH halves of
admissibility. The closed forms below follow from the von Mangoldt recursion
$a_n \log n = \sum_{d \mid n} \Lambda_{\mathrm{DH}}(d)\, a_{n/d}$ on the
periodic coefficient array and match the stored record digit-for-digit at the
printed precision:

\begin{center}
\small
\begin{tabular}{@{}c l l >{\raggedright\arraybackslash}p{4.1cm}@{}}
\toprule
$n$ & exact form & value & reading \\
\midrule
$3$ & $-\kappa \cdot \log 3$ & $\mathbf{-0.3120927\ldots}$ &
negativity at a \textbf{prime} (cheapest witness) \\
\addlinespace
$4$ & $-(2 + \kappa^2) \cdot \log 2$ & $\mathbf{-1.44223196\ldots}$ &
negativity at a \textbf{prime power} \\
\addlinespace
$6$ & $+(1 + \kappa^2) \cdot \log 6$ & $\mathbf{+1.93635607662\ldots}$ &
support \textbf{off prime powers} ($6 = 2 \cdot 3$) \\
\addlinespace
$12$ & $-\kappa \cdot [\kappa^2 \log 6 + (2 + \kappa^2) \log 2 + \log 3]$ &
$\mathbf{-0.762877471988\ldots}$ &
negativity at a composite (the original one-line witness) \\
& $\quad = -\kappa (1 + \kappa^2) \cdot \log 12$ & & \\
\bottomrule
\end{tabular}
\end{center}

\medskip
$\Lambda_{\mathrm{DH}}(3) < 0$ and $\Lambda_{\mathrm{DH}}(4) < 0$ are the
adjudication-adopted ``cheaper witnesses'': negativity is visible on prime
powers alone, with no support-off-prime-powers preamble needed.

\subsection{Epstein, \texorpdfstring{$Q = x^2 + 5y^2$ ($D = -20$, $h = 2$)}{Q
= x\^{}2 + 5y\^{}2 (D = -20, h = 2)}; RH-false class, and
inadmissible}\label{sec:epstein}

``RH-false class'' is classical: for binary quadratic forms of class number
$h(D) > 1$, the Epstein zeta function has infinitely many zeros in
$\operatorname{Re} s > 1$ (Davenport--Heilbronn \cite{DH1936}) --- cited for
the label only; nothing in the exclusion computation below depends on it.

Normalization (as stored): $F(s) = \zeta_Q(s)/2$, $b_n = r_Q(n)/2$,
$b_1 = 1$; $\Lambda_Q$ from $-F'/F$ (independent of the constant). Exact
rational-log arithmetic, not floating point; numeric rendering at $40$ digits.
The two witnesses:

\begin{itemize}
\item Support off prime powers ($6 = 2 \cdot 3$): the
composition-law/primitivity conclusion fails ---
\begin{align*}
\Lambda_Q(6) &= 2 \log 2 + 2 \log 3 = 2 \log 6 \\
&= \mathbf{+3.583518938456110001624954716761404545446}.
\end{align*}
\item $\Lambda_Q < 0$: the effectivity/nef-positivity conclusion fails ---
\begin{align*}
\Lambda_Q(36) &= -4 \log 2 - 4 \log 3 = -4 \log 6 \\
&= \mathbf{-7.167037876912220003249909433522809090892}.
\end{align*}
\end{itemize}

Support off prime powers for $n \le 100$ at
$\{6, 14, 21, 36, 46, 54, 69, 84, 86, 94\}$; negative values at
$\{36, 54, 84\}$.

\medskip
\noindent\textbf{The two structural checks} (both passed exactly for all
$n \le 200$):

\begin{enumerate}
\item \textbf{Genus identity.}
$(r_{Q_1}(n) + r_{Q_2}(n))/2 = \sum_{d \mid n} \chi_{-20}(d)$ exactly, where
$Q_1 = x^2 + 5y^2$ (principal) and $Q_2 = 2x^2 + 2xy + 3y^2$ (non-principal)
--- the coefficient machinery is verified against class-field theory, not
merely enumerated.
\item \textbf{Contrast object.} $\zeta_K = \zeta \cdot L(\chi_{-20})$ (the
full class-group average, Euler product intact) has $\Lambda_K$ supported on
prime powers and $\ge 0$ over the same range. What is \textbf{proved} is the
pair: (a) $\zeta_Q$ fails admissibility (the exact witnesses above); (b) the
class-group average $\zeta_K$ --- a theorem-level Euler product, being a
Dedekind zeta --- satisfies it. The violation therefore \textbf{tracks exactly
the class-orbit-breaking}: the orbit average passes, the single-class
evaluation fails; evaluating a single CM point of the Eisenstein torus period
instead of the full class-group orbit is exactly what destroys Euler
factorization. That mechanism attribution is the interpretation of (a) $+$ (b)
--- the mechanism claim of the original Lemma~B, surviving here at witness
level.
\end{enumerate}

\subsection{The over-breadth witness: \texorpdfstring{$\chi_4$}{chi4}
(RH-true, and inadmissible)}\label{sec:chi4}

$\Lambda_{\chi_4}(3) = a_3 \cdot \log 3 = \mathbf{-\log 3 = -1.0986\ldots}
< 0$ (computed in the adjudication to $50$ digits). An RH-true, degree-one,
Euler-intact Dirichlet $L$-function fails the admissibility condition. The
filter is therefore \textbf{over-broad in a revealing way}: it is a
\emph{positivity} filter, not an Euler-product filter. Any $\mathrm{GL}(2)$
object with some $a_p < 0$ is excluded alongside DH.

\section{The Beurling-instantiability and degree-one-rigidity
clauses}\label{sec:clauses}

\paragraph{Beurling/DMV instantiability.} The axiom class as printed does
\textbf{NOT} exclude Beurling/DMV worlds. Beurling generalized number systems
(DMV: Diamond--Montgomery--Vorhauer \cite{DMV}; the Broucke-school factories
\cite{BDR, Br, BV}) carry a full Euler product, multiplicativity at every
generalized prime, and $\Lambda_P(n) \ge 0$ pointwise \textbf{automatically}
--- while violating RH maximally. Every \textbf{operative} coefficient-level
consequence of the axiom class (prime-power support $+$ $\Lambda \ge 0$) holds
in these worlds by construction; (PF5), an uninstantiated template, contributes
no operative coefficient-level consequence to check --- its exclusion from this
accounting, and why imposing its FE shadow by hand would change nothing, are
stated in the PF5 status paragraph, \S\ref{sec:axioms}. \textbf{The axiom a
DMV world violates: currently NONE.} The candidate DMV-violated input, named
per the adjudicated repair, is the substrate-level theta-effectivity
$h^0_\theta(D) \ge \deg D + O(1)$ (Bost \cite{Bost}; van der Geer--Schoof
\cite{vdGS}) --- a substrate-level input no current object supplies. Until
such an axiom exists and is instantiated, the class cannot decide RH.

\paragraph{What the filter is.} The polarized-Frobenius axiom class as printed
is a \textbf{POSITIVITY filter separating $\{\mathrm{DH}, \mathrm{Epstein}\}$
from $\{\zeta\} \cup \{\text{all Beurling worlds}\}$}, and it additionally
\textbf{excludes RH-true signed-coefficient $L$-functions}
($\Lambda_{\chi_4}(3) = -\log 3 < 0$, \S\ref{sec:chi4}). At the coefficient
level the substrate could only ever attach to nonnegative-$\Lambda$ objects
($\zeta$, $\zeta_K$, Rankin--Selberg squares). It is not an Euler-product
filter and not an RH filter.

\paragraph{Degree-one rigidity: the closing bracket.} The two-sided
calibration any de-novo axiom set must clear: below $\{$FE $+$ Euler product
$+$ Ramanujan$\}$ at degree one, counterexamples are mass-produced (the
Beurling factory above); at them, rigidity --- Selberg class degree $1$ $=$
$\zeta$ $+$ shifted Dirichlet $L$-functions only, no elements of degree
$1 < d < 2$ (\textbf{Kaczorowski--Perelli} \cite{KP1, KP7}). The printed
axioms do not derive this classification; it is hereby cited as the
\textbf{closing-bracket target} the axiom class must eventually contain. The
gap between Beurling-factory death and Kaczorowski--Perelli rigidity is where
a sufficient axiom set must sit; the present class sits below it.

\paragraph{Relocation of RH content.} Consequently, \textbf{all RH content
must sit in the substrate $+$ Hodge-index composite} (the
substrate-and-positivity chain of the surrounding program; see the Provenance
remark for its internal names). This note carries credit ONLY in the axiom-format role: it
names the axioms a future substrate must instantiate and records which objects
provably cannot satisfy their coefficient-level consequences. Arguments
consuming ``multiplicativity as geometry'' should cite this note for the axiom
statement and the \textbf{Blomer--Leung beyond-endoscopy monoid converse}
\cite{BL} as the sharper formal home of the same exclusion --- one axiom, two
dialects.

\paragraph{Consistency rider.} This section implements, in writing, the
standing rider on the underlying barrier records: any future axiom-level
filter of this kind must add FE/duality $+$ degree-one-rigidity clauses and
relocate RH content to the substrate $+$ generator composite.

\section{What this note may NOT be cited for (adjudicated
exclusions)}\label{sec:notcite}

\begin{enumerate}
\item \textbf{Not ``provable today.''} Struck by the adjudication; the phrase
may not be revived for any part of this material beyond the
\S\ref{sec:admissibility} monoid unpacking, which is quasi-definitional.
\item \textbf{Not a theorem.} ``DH/Epstein admit no polarized Frobenius
system'' is not a theorem under either reading of the definition
(\S\ref{sec:fork}); the citable statement is inadmissibility
(\S\S\ref{sec:admissibility}--\ref{sec:witnesses}).
\item \sloppy \textbf{Not a discharge of the $\mathbb{F}_1$-model question.}
The question asks for a proof that DH/Epstein have no \emph{model} over
$\mathbb{F}_1$; that requires a formalized notion of $\mathbb{F}_1$-model and
remains open.
\item \textbf{No credit beyond the axiom-format role.} In particular: no claim
that the axiom class excludes Beurling/DMV worlds (\S\ref{sec:clauses} states
the opposite); no vacuous ``zeta instantiates the axioms'' credit (the
$\mathbb{P}^1 \times \mathbb{P}^1$ system of \S\ref{sec:hornA} shows
axiom-level instantiation can be worthless); no full credit while no zeta
instantiation of the axioms exists at any rank.
\item \textbf{Not a revival of per-prime-fiber assemblies.} Nothing here bears
on per-prime-fiber assemblies (products of elliptic curves attached to
distinct primes); those constructions are excluded on independent grounds
recorded in the surrounding program and remain excluded.
\end{enumerate}

\begin{provenance}
\sloppy
This note originated as deliverable M0 (axiom-format note) of the C3-r
direction of an internal research program on the Riemann hypothesis, written in
Session~6, 2026-08-26, per a binding adjudication of 2026-08-26. In the
program's terminology: the ``killers'' are two independent adversarial
reviewers who refuted the original Lemma~B theorem framing by blind
replication (\texttt{results/verdicts-c3d1.json}, kill:C3 fatal~2, kill2:C3
fatal~1); the ``adjudication'' is the binding third pass that re-derived the
findings and ordered the repairs (\texttt{results/adjudication-C3.json}, the
mandatory repair beginning ``Reprice M0 \ldots''; adjudication entries are
cited by quoted opening phrase, never by array index); the recommission order
``C3-r'' is in \texttt{directions/}\allowbreak
\texttt{C3-geometric-substrate.md}. The consistency
obligations engaged in the body correspond to the program's barrier-catalog
entries I.1 (the Blomer--Leung monoid converse as the sharper formal home),
I.2 (the Beurling/DMV factory), I.7 (degree-one rigidity), III.21, and IV.10
(the rider implemented by \S\ref{sec:clauses}; the same entry records the
exclusion of the per-prime-fiber assemblies of \S\ref{sec:notcite},
item~5) in \texttt{BARRIER-ZOO.md}. The ``missing substrate'' of
\S\ref{sec:hornB} and \S\ref{sec:admissibility} is the program's item M2; the
``substrate $+$ Hodge-index composite'' of \S\ref{sec:clauses} is its M1 gate
$\to$ M2c substrate $\to$ M3/M4 positivity chain, and the theta-effectivity
axiom named there is its M4-level input.

Numeric traceability (paths relative to the program root
\texttt{anthropic/rh-program/}): the fork, the
$\mathbb{P}^1 \times \mathbb{P}^1$ counterexample, bidegree multiplicativity,
$(\Gamma_n \cdot \Delta) = n+1$, the non-multiplicativity $7 \ne 12$, the
char-$p$ correction, and the $\mathbb{F}_7$ verification are in
\texttt{results/verdicts-c3d1.json} and \texttt{results/adjudication-C3.json}
(computation~b); $\kappa$ (41 digits) and
$\Lambda_{\mathrm{DH}}(3), (4), (6), (12)$ in adjudication computation~d and
three independent computations, digit-for-digit, with the underlying DH
machinery in \texttt{results/decisive-tests/ccm-dh-filter.json} and
\texttt{results/ccm-dh-test/}; the Epstein exact pair, 40-digit values,
support/negativity lists, genus identity, and $\zeta_K$ contrast in
\texttt{results/c3-m0-epstein/}\allowbreak
\texttt{n1\_epstein\_witness.json} (with its generating
script); $\Lambda_{\chi_4}(3) = -\log 3$ in adjudication computation~c. The
exact expressions $-\kappa \log 3$, $-(2+\kappa^2)\log 2$,
$+(1+\kappa^2)\log 6$, and
$-\kappa[\kappa^2 \log 6 + (2+\kappa^2)\log 2 + \log 3]$ were unpacked from
the standard von Mangoldt recursion on the periodic DH coefficient array; each
was checked against the stored decimal record and agrees to all printed
digits. The Epstein closed forms are stored verbatim. An independent referee
pass of 2026-08-26 (\texttt{results/c3-r/referee-m0.md}) returned PASS WITH
REPAIRS --- $0$ fatal, $3$ major, $7$ minor; no mathematical error; all six
witnesses, both closed-form families, the $\mathbb{F}_7$ accounting, and the
$\mathbb{P}^1 \times \mathbb{P}^1$ counterexample were replicated from
scratch, and the DH functional equation and off-line zero were independently
re-verified to residuals $\le 10^{-60}$; all repairs were applied the same
day, none rejected or deferred.
\end{provenance}

\begin{thebibliography}{99}

\bibitem{BL}
V.~Blomer and W.~H.~Leung,
\emph{A GL(3) converse theorem via a ``beyond endoscopy'' approach},
Adv. Math. \textbf{485} (2026), Art. 110716;
arXiv:2401.04037.

\bibitem{Bost}
J.-B.~Bost,
\emph{Theta invariants of Euclidean lattices and infinite-dimensional
Hermitian vector bundles over arithmetic curves},
arXiv:1512.08946.

\bibitem{Br}
F.~Broucke,
\emph{On the connection between zero-free regions and the error term in the
Prime Number Theorem},
arXiv:2507.13780.
(Cited for its construction of Beurling zeta functions with infinitely many
zeros on a prescribed contour and none to the right of it.)

\bibitem{BDR}
F.~Broucke, G.~Debruyne, and Sz.~Gy.~R\'ev\'esz,
\emph{Some examples of well-behaved Beurling number systems},
Trans. Amer. Math. Soc. \textbf{378} (2025), 477--501;
arXiv:2309.01567.
(Cited for its $[\alpha, \beta]$-systems of Beurling generalized primes.)

\bibitem{BV}
F.~Broucke and J.~Vindas,
\emph{A new generalized prime random approximation procedure and some of its
applications},
Math. Z. \textbf{307} (2024), Art. 62;
arXiv:2102.08478.
(Cited for its discretization of continuous Beurling systems preserving
$\Lambda_P \ge 0$.)

\bibitem{CM}
H.~Chen and A.~Moriwaki,
\emph{Arakelov Geometry over Adelic Curves},
Lecture Notes in Mathematics 2258, Springer, 2020;
arXiv:1903.10798.

\bibitem{DH1936}
H.~Davenport and H.~Heilbronn,
\emph{On the zeros of certain Dirichlet series (Second paper)},
J. London Math. Soc. \textbf{11} (1936), 307--312.

\bibitem{DMV}
H.~G.~Diamond, H.~L.~Montgomery, and U.~M.~A.~Vorhauer,
\emph{Beurling primes with large oscillation},
Math. Ann. \textbf{334} (2006), 1--36.

\bibitem{vdGS}
G.~van der Geer and R.~Schoof,
\emph{Effectivity of Arakelov divisors and the theta divisor of a number
field},
Selecta Math. (N.S.) \textbf{6} (2000), 379--398;
arXiv:math/9802121.

\bibitem{KP1}
J.~Kaczorowski and A.~Perelli,
\emph{On the structure of the Selberg class, I: $0 \le d \le 1$},
Acta Math. \textbf{182} (1999), 207--241.

\bibitem{KP7}
J.~Kaczorowski and A.~Perelli,
\emph{On the structure of the Selberg class, VII: $1 < d < 2$},
Ann. of Math. (2) \textbf{173} (2011), 1397--1441.

\bibitem{YZ}
X.~Yuan and S.-W.~Zhang,
\emph{Adelic line bundles on quasi-projective varieties},
Annals of Mathematics Studies 223, Princeton University Press;
arXiv:2105.13587.

\end{thebibliography}

\end{document}
